\section{Experiments}
\label{sec:Experiments}

\subsection{Experimental Settings}

\textbf{Datasets and evaluation protocol.} To pre-train ImageED, the ED-aware teacher, the ED predictor, and the reliability estimator, we select the first 2 million unlabeled molecular conformations and their DFT-computed ED data from the EDBench database~\cite{xiang2025edbenchlargescaleelectrondensity}, generated by Psi4 software~\cite{turney2012psi4} with a grid spacing of 0.4 Bohr. In the evaluation stage, we select 12 widely used tasks related to quantum mechanical properties from QM9~\cite{ramakrishnan2014quantum} and 10 common tasks related to energy/force from revised MD17 (rMD17)~\cite{christensen2020role}. For force prediction, we first predict the molecular energy and use the gradient of each node position as the force, \ie, \textit{force=-torch.autograd.grad(outputs=energy, inputs=positions)} in PyTorch~\cite{paszke2019pytorch}. The dataset split follows Geom3D~\cite{liu2024symmetry}: 110K for training, 10K for validation, and 11K for testing in QM9; 950 for training, 50 for validation, and 1000 for testing in rMD17. For the out-of-distribution case study (Sec.~\ref{sec:case_study_chignolin}), we additionally use the AIMD-Chig dataset~\cite{wang2023aimdchig}: 9{,}543 Chignolin (166-atom mini-protein) conformations sampled by replica-exchange MD and labeled at DFT, split 8:1:1 into 7634/954/955 for train/validation/test. We use mean absolute error (MAE) as the evaluation metric.

\textbf{Baselines.} To evaluate EDG and EDG++, we select geometry-based models with diverse architectures as geometry students. For QM9, we evaluate SchNet~\cite{schutt2017schnet}, Equiformer~\cite{liao2023equiformer}, and SphereNet~\cite{liu2022spherical}. For rMD17, we evaluate SchNet, SphereNet, and ViSNet~\cite{wang2024enhancing}, applying both EDG and EDG++ to all three architectures. Following~\cite{liu2022spherical,wang2024enhancing,liu2024symmetry}, we ensure that each baseline is fully trained with the same hyperparameters as in the original papers.

\textbf{Implementation details.} The implementation details of ImageED pre-training, ED-aware teacher pre-training, and reliability estimator training are provided in Appendix~\ref{sec:appendix_implementation}. In the distillation stage, we conduct grid search over the hyperparameters described in Appendix~\ref{sec:appendix_hyperparam}. The mapper $f_M$ and task predictor $f_T$ each consist of a single linear layer. We select the best checkpoint based on validation MAE for QM9 and validation force MAE for rMD17.

\textbf{Hyperparameter search and evaluation protocol.} We follow a strict protocol to avoid data leakage: hyperparameter search is performed exclusively on the validation set, and all reported results are evaluated on the held-out test set. Following standard practice in molecular property prediction~\cite{liu2022spherical,wang2024enhancing}, we report per-task-best results across hyperparameter configurations for all methods, ensuring a fair comparison on the same search budget.

\textbf{Single-seed protocol.} All experiments---including the Chignolin case study training---use a single random seed (seed=42); the 10 MD trajectories in the case study (Sec.~\ref{sec:case_study_chignolin}) are launched from 10 distinct test-set starting conformations with a single trained model per method, not from independent training seeds. We assess statistical reliability via: (1)~paired bootstrap CIs on per-sample test predictions (Appendix~\ref{sec:appendix_bootstrap}); (2)~sign tests (34/36 QM9 wins: $p < 10^{-8}$; 55/60 rMD17: $p < 10^{-11}$); and (3)~coefficient of variation across 45 configurations (Section~\ref{sec:model_robustness}). Multi-seed evaluation remains a limitation (Section~\ref{sec:Conclusion}).

\subsection{Main Results on QM9}

\begin{table*}[!tb]
  \scriptsize
  \centering
  \caption{MAE performance on 12 quantum mechanical property prediction tasks in QM9. ``Baseline'' denotes the original geometry model without distillation. ``EDG'' denotes naive distillation (conference version). ``EDG++'' denotes reliability-aware selective distillation; for each property, we report the best result across the hyperparameter search. $\Delta$ represents the relative improvement of EDG++ over the baseline. Bold indicates the best result in each column. $\downarrow$ indicates lower is better.}
  \label{tab:qm9_edgpp}
    \resizebox{\textwidth}{!}{%
    \begin{tabular}{clcccccccccccc}
    \toprule
        & \multirow{2}{*}{Method} & $\alpha$ $\downarrow$ & $\Delta\varepsilon$ $\downarrow$ & $\varepsilon_{\text{HOMO}}$ $\downarrow$ & $\varepsilon_{\text{LUMO}}$ $\downarrow$ & $\mu$ $\downarrow$ & $C_v$ $\downarrow$ & $G$ $\downarrow$ & $H$ $\downarrow$ & $R^2$ $\downarrow$ & $U$ $\downarrow$ & $U_0$ $\downarrow$ & ZPVE $\downarrow$\\
        & & $a_0^3$ & $meV$ & $meV$ & $meV$ & $D$ & $\frac{cal}{mol\cdot K}$ & $meV$ & $meV$ & $a_0^2$ & $meV$ & $meV$ & $meV$\\
    \midrule
    \multirow{4}{*}{\rotatebox{90}{SchNet}}
    & Baseline & 0.07021 & 50.829 & 31.952 & 26.168 & 0.03013 & 0.03228 & 14.678 & 14.090 & 0.13455 & 14.142 & 13.915 & 1.714 \\
    & EDG & 0.06866 & 49.778 & 31.884 & 25.972 & \textbf{0.02980} & 0.03162 & 14.022 & 13.841 & \textbf{0.12458} & 13.794 & 13.826 & 1.688 \\
    & EDG++ & \textbf{0.06810} & \textbf{49.426} & \textbf{31.385} & \textbf{25.499} & 0.02984 & \textbf{0.03147} & \textbf{13.729} & \textbf{13.773} & 0.12591 & \textbf{13.613} & \textbf{13.416} & \textbf{1.674} \\
    & $\Delta$ & $\uparrow$3.0\% & $\uparrow$2.8\% & $\uparrow$1.8\% & $\uparrow$2.6\% & $\uparrow$1.0\% & $\uparrow$2.5\% & $\uparrow$6.5\% & $\uparrow$2.2\% & $\uparrow$6.4\% & $\uparrow$3.7\% & $\uparrow$3.6\% & $\uparrow$2.3\% \\
    \midrule
    \multirow{4}{*}{\rotatebox{90}{\scriptsize Equiformer}}
    & Baseline & 0.06762 & 46.308 & 26.017 & 23.681 & 0.02074 & 0.02733 & 18.439 & 16.453 & 0.45828 & 15.339 & 23.928 & 1.537 \\
    & EDG & 0.06476 & 45.813 & 25.492 & 23.266 & \textbf{0.01985} & 0.02642 & 15.976 & \textbf{14.451} & 0.43947 & 15.466 & \textbf{16.517} & \textbf{1.529} \\
    & EDG++ & \textbf{0.06175} & \textbf{45.089} & \textbf{25.269} & \textbf{22.283} & 0.02113 & \textbf{0.02637} & \textbf{12.058} & 14.709 & \textbf{0.40354} & \textbf{14.802} & 16.918 & 1.586 \\
    & $\Delta$ & $\uparrow$8.7\% & $\uparrow$2.6\% & $\uparrow$2.9\% & $\uparrow$5.9\% & $\downarrow$1.9\% & $\uparrow$3.5\% & $\uparrow$34.6\% & $\uparrow$10.6\% & $\uparrow$11.9\% & $\uparrow$3.5\% & $\uparrow$29.3\% & $\downarrow$3.2\% \\
    \midrule
    \multirow{4}{*}{\rotatebox{90}{SphereNet}}
    & Baseline & 0.04670 & 40.129 & 22.007 & 19.435 & 0.02689 & 0.02437 & 7.875 & 7.199 & 0.25821 & 6.999 & 6.641 & 1.253 \\
    & EDG & 0.04592 & 39.694 & 21.842 & 19.014 & \textbf{0.02648} & \textbf{0.02376} & 7.769 & \textbf{6.283} & \textbf{0.24935} & 6.502 & \textbf{6.101} & 1.206 \\
    & EDG++ & \textbf{0.04567} & \textbf{39.470} & \textbf{21.766} & \textbf{18.750} & 0.02652 & 0.02400 & \textbf{7.656} & 6.364 & 0.25341 & \textbf{6.409} & 6.424 & \textbf{1.198} \\
    & $\Delta$ & $\uparrow$2.2\% & $\uparrow$1.6\% & $\uparrow$1.1\% & $\uparrow$3.5\% & $\uparrow$1.4\% & $\uparrow$1.5\% & $\uparrow$2.8\% & $\uparrow$11.6\% & $\uparrow$1.9\% & $\uparrow$8.4\% & $\uparrow$3.3\% & $\uparrow$4.4\% \\
    \midrule
    \multicolumn{14}{l}{\scriptsize\textit{Avg./median $\Delta$: SchNet $\uparrow$3.2\%/2.7\%, Equiformer $\uparrow$9.0\%/4.7\%, SphereNet $\uparrow$3.6\%/2.5\%. Overall 34/36 tasks improved. Equiformer's high mean is partly driven by thermodynamic outliers ($G$ +34.6\%, $U_0$ +29.3\%); see Sec.~IV-A property-dependent analysis.}}\\
    \bottomrule
    \end{tabular}%
    }
\end{table*}

Table~\ref{tab:qm9_edgpp} presents the main QM9 results. EDG++ improves 34 out of 36 task-model combinations (sign test $p < 10^{-8}$), with average relative gains of 3.2\% (SchNet), 9.0\% (Equiformer), and 3.6\% (SphereNet). Comparing EDG++ with EDG, selective distillation outperforms naive distillation on the majority of tasks (SchNet: 10/12, SphereNet: 7/12), confirming that reliability-aware filtering successfully excludes unreliable teacher predictions.

\textbf{Property-dependent improvements.} The improvements exhibit a physically interpretable hierarchy consistent with the Hohenberg-Kohn theorem: thermodynamic properties achieve the largest and most consistent improvements (5/5 win rate across all architectures; avg.\ +3.70\% to +16.30\%), as these ground-state observables are uniquely determined by the electron density~\cite{hohenberg1964inhomogeneous}. The only two degradation cases---Equiformer on $\mu$ ($\downarrow$1.9\%) and ZPVE ($\downarrow$3.2\%)---are architecture-specific and improve on SchNet and SphereNet, confirming this is not a fundamental limitation. Figure~\ref{fig:property_group} visualizes this hierarchy across all 36 combinations; the detailed property-group breakdown is in Appendix~\ref{sec:appendix_property_group}.

\begin{figure*}[!t]
\centering
\includegraphics[width=0.95\textwidth]{imgs_edgpp/experiments/property_group_improvement.pdf}
\caption{Relative improvement (\%) for all 36 QM9 task-model combinations (3 architectures $\times$ 12 properties). Properties are grouped by Hohenberg-Kohn relevance: thermodynamic (left), electronic (middle), and other (right). Blue cells indicate improvement; warm-colored cells indicate degradation.}
\label{fig:property_group}
\end{figure*}

\subsection{Tail Error and Robustness Analysis}
\label{sec:error_distribution}

To test the heavy-tailed error structure assumed in Proposition~\ref{prop:gradient_noise}, we analyze how threshold tuning affects prediction errors across difficulty regimes. We compare two EDG++ configurations---tuned $\kappa$ versus default $\kappa{=}0$---on four representative tasks spanning different architectures and property types.

\begin{figure*}[!t]
\centering
\includegraphics[width=0.95\textwidth]{imgs_edgpp/experiments/tail_error_ecdf.pdf}
\caption{Cumulative mean absolute error curves on QM9 for four representative tasks. The x-axis (log scale) shows the top $K$\% of samples ranked by difficulty (default-threshold error magnitude), from the hardest 0.5\% to 25\%. The EDG++ tuned-$\kappa$ curve (orange) lies consistently below the EDG++ default-threshold ($\kappa\!=\!0$) curve (blue). Annotated percentages show the relative error reduction at the 1\% and 5\% marks. \textit{Note}: both curves use the reliability estimator; the comparison isolates the effect of threshold tuning. Difficulty ranking is defined by the default-threshold model's errors; y-axis scales differ across panels to maximize readability within each task.}
\label{fig:tail_error}
\end{figure*}

Figure~\ref{fig:tail_error} shows that threshold tuning substantially reduces tail errors: at the top 1\%, reductions range from 27.0\% (SphereNet-$H$) to 37.9\% (Equiformer-$U_0$), providing indirect evidence for the heavy-tailed error structure assumed in Proposition~\ref{prop:gradient_noise} (see Section~\ref{sec:Theoretical_Analysis} for the causal link between downstream prediction errors and teacher reconstruction errors). The tuned curve remains below the default curve across the full visible range (top 0.5\%--25\%), indicating that the benefit extends broadly across the harder half of the test set. Standard deviation reductions of 5.3--17.3\% and maximum error reductions up to 39.1\% (Equiformer-$U_0$) confirm improved robustness. Difficulty-stratified analysis (Appendix~\ref{sec:appendix_difficulty_stratification}) reveals a capacity reallocation effect: threshold tuning shifts learning capacity from easy samples (where the baseline already performs well) to hard samples (where stricter filtering provides the most corrective signal), consistent with Proposition~\ref{prop:gradient_noise}.


\subsection{Main Results on rMD17}

\begin{table*}[!tb]
  \scriptsize
  \centering
  \caption{MAE on 10 energy ($\frac{kcal}{mol}$) prediction tasks in rMD17. Baseline and EDG results are from the conference version~\cite{xiang2025edg}. Bold indicates the best result per column within each architecture block. Hyperparameter configurations, search space, and progressive-ablation study are provided in Appendix~\ref{app:hyper}--\ref{app:progressive}. Malona.\ and Naphth.\ abbreviate Malonaldehyde and Naphthalene.}
  \label{tab:rmd17_energy}
  \setlength{\tabcolsep}{4pt}
    \begin{tabular}{clcccccccccc}
    \toprule
    & Method & Aspirin $\downarrow$ & Azobenz. $\downarrow$ & Benzene $\downarrow$ & Ethanol $\downarrow$ & Malona. $\downarrow$ & Naphth. $\downarrow$ & Paracet. $\downarrow$ & Salicylic $\downarrow$ & Toluene $\downarrow$ & Uracil $\downarrow$ \\
    \midrule
    \multirow{3}{*}{\tabincell{c}{SchNet}}
    & Baseline & 0.73909 & 0.39678 & 0.02052 & 0.12516 & 0.16142 & 0.21158 & 0.37097 & 0.19078 & 0.20797 & 0.07872 \\
    & EDG & 0.35525 & 0.33441 & 0.01711 & 0.06061 & 0.11181 & 0.07338 & 0.28303 & 0.15697 & 0.08780 & \textbf{0.07433} \\
    & EDG++ & \textbf{0.30625} & \textbf{0.30013} & \textbf{0.01683} & \textbf{0.06056} & \textbf{0.10733} & \textbf{0.07069} & \textbf{0.19075} & \textbf{0.15447} & \textbf{0.08319} & 0.07598 \\
    \midrule
    \multirow{3}{*}{\tabincell{c}{SphereNet}}
    & Baseline & 0.18091 & 0.09794 & 0.00647 & 0.03784 & 0.06005 & 0.03823 & 0.10425 & 0.14119 & 0.03452 & 0.08088 \\
    & EDG & 0.13622 & \textbf{0.06788} & 0.00413 & \textbf{0.03575} & \textbf{0.05659} & 0.02753 & 0.09934 & 0.09569 & 0.02413 & \textbf{0.03683} \\
    & EDG++ & \textbf{0.12391} & 0.07125 & \textbf{0.00341} & \textbf{0.03575} & 0.06611 & \textbf{0.02256} & \textbf{0.09628} & \textbf{0.06563} & \textbf{0.02220} & 0.04841 \\
    \midrule
    \multirow{3}{*}{\tabincell{c}{ViSNet}}
    & Baseline & 0.05547 & 0.02081 & 0.00627 & 0.01095 & 0.01517 & 0.01313 & 0.02700 & 0.01966 & 0.01089 & 0.01238 \\
    & EDG & 0.04650 & 0.01838 & 0.00616 & 0.00990 & 0.01395 & 0.01178 & 0.02491 & 0.01906 & 0.00998 & \textbf{0.01188} \\
    & EDG++ & \textbf{0.04541} & \textbf{0.01588} & \textbf{0.00572} & \textbf{0.00958} & \textbf{0.01367} & \textbf{0.00930} & \textbf{0.02369} & \textbf{0.01547} & \textbf{0.00798} & 0.01192 \\
    \bottomrule
    \end{tabular}
\end{table*}

\begin{table*}[!tb]
  \scriptsize
  \centering
  \caption{MAE performance on 10 force ($\frac{kcal}{mol\cdot \text{\r{A}}}$) prediction tasks in rMD17. Setup follows Table~\ref{tab:rmd17_energy}.}
  \label{tab:rmd17_force}
  \setlength{\tabcolsep}{4pt}
    \begin{tabular}{clcccccccccc}
    \toprule
    & Method & Aspirin $\downarrow$ & Azobenz. $\downarrow$ & Benzene $\downarrow$ & Ethanol $\downarrow$ & Malona. $\downarrow$ & Naphth. $\downarrow$ & Paracet. $\downarrow$ & Salicylic $\downarrow$ & Toluene $\downarrow$ & Uracil $\downarrow$ \\
    \midrule
    \multirow{3}{*}{\tabincell{c}{SchNet}}
    & Baseline & \textbf{1.04245} & 0.90082 & 0.18569 & 0.38519 & 0.65536 & 0.39851 & 0.82544 & 0.77487 & 0.48322 & 0.51399 \\
    & EDG & 1.04910 & 0.91692 & 0.17113 & 0.37901 & 0.64678 & 0.39509 & 0.83296 & \textbf{0.74308} & \textbf{0.47790} & \textbf{0.50783} \\
    & EDG++ & 1.05170 & \textbf{0.89131} & \textbf{0.16655} & \textbf{0.37584} & \textbf{0.63879} & \textbf{0.38576} & \textbf{0.77857} & 0.74746 & 0.47876 & 0.54119 \\
    \midrule
    \multirow{3}{*}{\tabincell{c}{SphereNet}}
    & Baseline & 0.39134 & 0.21776 & 0.02151 & 0.19432 & 0.29278 & 0.11141 & 0.32265 & 0.28692 & 0.10978 & 0.27702 \\
    & EDG & 0.38598 & \textbf{0.21665} & 0.02101 & 0.18741 & \textbf{0.28456} & 0.11102 & 0.32023 & \textbf{0.28299} & 0.10804 & \textbf{0.17179} \\
    & EDG++ & \textbf{0.37929} & 0.21708 & \textbf{0.02012} & \textbf{0.18698} & 0.31599 & \textbf{0.09541} & \textbf{0.31846} & 0.29452 & \textbf{0.10341} & 0.22544 \\
    \midrule
    \multirow{3}{*}{\tabincell{c}{ViSNet}}
    & Baseline & 0.15164 & 0.05729 & 0.00656 & 0.05688 & 0.09275 & 0.02808 & 0.10488 & 0.08348 & 0.02980 & 0.05252 \\
    & EDG & 0.14996 & \textbf{0.05691} & 0.00647 & 0.05558 & 0.08992 & 0.02798 & 0.10599 & 0.08107 & 0.02780 & 0.05100 \\
    & EDG++ & \textbf{0.14813} & 0.05712 & \textbf{0.00619} & \textbf{0.05508} & \textbf{0.08909} & \textbf{0.02713} & \textbf{0.10410} & \textbf{0.07959} & \textbf{0.02672} & \textbf{0.05085} \\
    \bottomrule
    \end{tabular}
\end{table*}

Tables~\ref{tab:rmd17_energy} and~\ref{tab:rmd17_force} present the rMD17 results. EDG delivers consistent energy improvements across all architectures (33.7\% for SchNet and 25.2\% for SphereNet on average), while force improvements remain modest (1.5--5.3\%), because forces depend on local atomic interactions rather than the global energy distribution captured by ED. EDG++ provides further gains on top of EDG across all three architectures. Aggregated across the 60 molecule-task combinations (30 energy + 30 force), EDG++ improves over the no-distillation baseline on 55/60 tasks. Hyperparameter configurations, a progressive-component ablation isolating the contribution of the reliability estimator, and per-molecule error analyses (including a boundary case on malonaldehyde) are provided in Appendix~\ref{app:hyper}--\ref{app:boundary}.

\FloatBarrier
\subsection{Ablation Studies}

\subsubsection{Selective Distillation Components}

\begin{table}[!t]
  \scriptsize
  \centering
  \caption{Ablation on thresholding strategy ($\beta$) for QM9. We report the number of tasks (out of 12) where each $\beta$ value achieves the best result. ``Total'' aggregates across all three architectures (out of 36).}
  \label{tab:ablation_beta}
    \begin{tabular}{lccccc}
    \toprule
    \textbf{Strategy} & \textbf{$\beta$} & \textbf{SchNet} & \textbf{SphereNet} & \textbf{Equiformer} & \textbf{Total} \\
    \midrule
    Global only & 0 & 2 & 3 & 7 & 12/36 \\
    Global+Local & 0.5 & 5 & 7 & 5 & \textbf{17/36} \\
    Local only & 1.0 & 5 & 2 & 0 & 7/36 \\
    \bottomrule
    \end{tabular}
\end{table}

Table~\ref{tab:ablation_beta} analyzes the mixing coefficient $\beta$. The hybrid strategy ($\beta=0.5$) achieves the most wins overall (17/36). The preference is architecture-dependent: Equiformer favors global thresholding (7/12), while SchNet and SphereNet benefit from hybrid or local strategies, suggesting that architectures with global attention already capture batch-level variation internally.

\begin{table}[!t]
  \scriptsize
  \centering
  \caption{Incremental contribution of threshold tuning on QM9. ``$\kappa\!=\!0$'' denotes the default-threshold configuration. $\Delta_{\text{inc}}$ is the incremental improvement from tuning $\kappa$ beyond the default.}
  \label{tab:incremental}
    \begin{tabular}{lccc}
    \toprule
    \textbf{Model} & \textbf{Improved} & \textbf{Avg.\ $\Delta_{\text{inc}}$} & \textbf{Max $\Delta_{\text{inc}}$} \\
    \midrule
    SchNet & 10/12 & +0.97\% & +2.46\% ($R^2$) \\
    SphereNet & 9/12 & +1.00\% & +3.00\% (ZPVE) \\
    Equiformer & 8/12 & +2.87\% & +12.82\% ($H$) \\
    \bottomrule
    \end{tabular}
\end{table}

Table~\ref{tab:incremental} isolates the contribution of threshold tuning by comparing tuned $\kappa$ against the default $\kappa{=}0$ configuration. Even with $\kappa=0$, the reliability estimator is active and filters samples below the mean reliability score; tuning $\kappa$ provides consistent incremental improvements across architectures (avg.\ +0.97\% to +2.87\%). Equiformer benefits the most, consistent with its higher sensitivity to filtering thresholds. % [Polish] Original: The consistent improvement across 27--31 out of 36 tasks provides indirect but strong evidence for the reliability estimator's effectiveness...
The fact that 27--31 out of 36 tasks improve provides indirect but strong evidence for the reliability estimator: if the reliability ranking were uncorrelated with actual teacher error, threshold-based filtering would not produce systematic gains.

\subsubsection{Hyperparameter Sensitivity}
\label{sec:model_robustness}

The distillation weight $\lambda$ exhibits architecture-dependent sensitivity (detailed per-task results in Appendix~\ref{sec:appendix_lambda}): Equiformer averages 7.77\% sensitivity (5$\times$ SchNet's 1.47\%), with thermodynamic properties most sensitive ($U_0$: 34.87\%). SchNet maintains consistently low sensitivity (max 2.37\%). The coefficient of variation (CV) of MAE across all 45 configurations further confirms robustness: SchNet (mean CV = 1.1\%), SphereNet (1.8\%), and Equiformer (4.21\%). Additional hyperparameter interaction analysis is in Appendix~\ref{sec:appendix_heatmap}.

\FloatBarrier

% ============================================================================
% CHIGNOLIN CASE STUDY
% ============================================================================
\subsection{Case Study: Observing Negative Transfer in the Wild (Chignolin, 166-atom Mini-Protein)}
\label{sec:case_study_chignolin}

% [Smoking-gun framing] Position the case study as the empirical observation that
% the negative-transfer mechanism EDG++ defends against truly manifests at scale.
\textbf{From theory to a stress test.} Proposition~\ref{prop:gradient_noise} predicts that selective distillation reduces gradient noise whenever teacher reliability and error are positively associated---a regime that intensifies as the teacher leaves its training envelope. The reliability estimator was therefore designed to gate noisy supervision at training time, but because it scores any input the teacher sees, the same frozen module can flag distributional drift at inference. Chignolin lets us test both roles simultaneously: a 166-atom mini-protein stresses the gating mechanism well outside the teacher's small-molecule training regime, and its long MD rollouts offer a natural setting to demonstrate the estimator as an inference-time OOD detector.

\textbf{Why this case study.} On QM9 and rMD17, EDG rarely hurts the baseline, so EDG++'s gating mostly outperforms EDG rather than rescuing it. To observe the negative-transfer failure mode that the theory anticipates---and to verify the gating mechanism recovers from it---we need a setting that pushes the teacher far enough out of its training distribution to expose the failure mode. This case study provides such a setting.

\textbf{Subject.} We choose Chignolin (CLN025 variant, sequence YYDPETGTWY), a 166-atom mini-protein whose backbone flexibility, aromatic side chains, and overall size are all substantially outside the PCQM4Mv2 conformations used to pre-train the ED teacher---an out-of-distribution regime far beyond the small-molecule benchmarks (QM9, rMD17; $<\!30$ atoms) on which EDG++ has so far been validated.

\textbf{Dataset and protocol.} We use the AIMD-Chig dataset~\cite{wang2023aimdchig}: 9{,}543 Chignolin conformations sampled by replica-exchange MD and labeled at DFT (B3LYP/6-31G(d,p)) with Gaussian 16~\cite{frisch2016gaussian16}, split 8:1:1 into train/validation/test. The student is ViSNet~\cite{wang2024enhancing}, chosen because it is the only architecture in our rMD17 table that scales to 166 atoms, and trained for 1{,}000 epochs with the same optimizer/lr schedule as the ViSNet paper's Chignolin configuration. We compare three variants: (i)~baseline ViSNet with no distillation; (ii)~ViSNet + EDG (naive distillation, fixed $\lambda$); and (iii)~ViSNet + EDG++ with the default-threshold configuration ($\kappa_{\text{all}}\!=\!0$, $\beta\!=\!0$), matching the out-of-the-box rMD17 protocol. The pre-trained ED-aware teacher and reliability estimator are frozen; only the geometry student is trained on Chignolin.

\textbf{Teacher reliability is domain-sensitive.} Before training, we pass the Chignolin structural images through the frozen teacher pipeline and compute reliability scores (negated evaluator output; higher = more reliable). Appendix~\ref{sec:appendix_chignolin_sanity} reports the full histogram against the 10 rMD17 molecules. The Chignolin distribution is shifted toward lower reliability (mean $-0.86$ vs.\ rMD17 mean $-0.68$ to $-0.79$, std $0.055$), with a tail overlapping the most reliable small-molecule regime. This is exactly the situation where naive distillation is most likely to hurt and where selective distillation should help: most Chignolin samples carry lower-quality teacher supervision, but a subset remains trustworthy.

\textbf{Per-conformation accuracy across the free-energy landscape.} Figure~\ref{fig:chignolin_case_study}(b) shows the 2D density landscape of the 9{,}543 DFT-labeled conformations along the ViSNet-paper reaction coordinates $d_{D3\text{-}G7}$ and $d_{E5\text{-}T8}$. The folded basin is concentrated on the left (low $d_{D3\text{-}G7}$), and the unfolded region spreads to the right. Six representative conformations are sampled from the held-out test set to illustrate the two distinct mechanisms quantified in the full test-set distribution (Table~\ref{tab:chignolin_energy_force}; rescue/amplify decomposition reported below): the \emph{rescue} mode, in which naive EDG degrades the baseline and EDG++ recovers (observed in 22.7\% of test conformations), and the \emph{amplify} mode, exhibiting strict monotone improvement (19.9\%). The 6 conformations are paired one rescue + one amplify per region (folded / transition / unfolded). EDG++ achieves the lowest absolute error on all 6 conformations; naive EDG degrades baseline on 3/6 (the rescue cases), matching the negative-transfer hypothesis. The effect is largest on unfolded-basin conformations, which exhibit the lowest teacher reliability and the largest prediction errors overall.

\textbf{Aggregate test-set accuracy.} On the Chignolin test set of 955 held-out conformations (Table~\ref{tab:chignolin_energy_force}), ViSNet+EDG++ reduces mean energy MAE by 6.8\% over baseline (1.2521 $\rightarrow$ 1.1671~kcal/mol), while naive EDG \emph{increases} mean energy MAE by 4.0\% (1.2521 $\rightarrow$ 1.3024~kcal/mol). Mean force MAE follows the identical monotonic pattern (baseline \SI{0.3075}{kcal/mol/\angstrom} $\to$ EDG \SI{0.3160}{kcal/mol/\angstrom} ($+\!2.8\%$) $\to$ EDG++ \SI{0.2985}{kcal/mol/\angstrom} ($-\!2.9\%$)). \textbf{This is the first systematic instance in our study of naive distillation \emph{measurably} degrading the baseline at scale, and the first instance of EDG++'s gating successfully \emph{recovering} from it---the central design hypothesis of the paper, observed in the wild.}

\begin{table}[!t]
  \scriptsize
  \centering
  \caption{Chignolin test-set accuracy on 955 held-out conformations. Energy in \SI{}{kcal/mol}, force in \SI{}{kcal/mol/\angstrom}. $\Delta$ relative to baseline.}
  \label{tab:chignolin_energy_force}
  \begin{tabular}{lcc}
    \toprule
    \textbf{Method} & \textbf{Energy MAE} $\downarrow$ & \textbf{Force MAE} $\downarrow$ \\
    \midrule
    ViSNet (baseline)  & 1.2521 & 0.3075 \\
    ViSNet + EDG (naive)   & 1.3024 ($\Delta\!+\!4.0\%$) & 0.3160 ($\Delta\!+\!2.8\%$) \\
    ViSNet + EDG++ ($\kappa{=}0$) & \textbf{1.1671} ($\Delta\!-\!6.8\%$) & \textbf{0.2985} ($\Delta\!-\!2.9\%$) \\
    \bottomrule
  \end{tabular}
\end{table}

\textbf{Stability under MD simulation.} We run 10 independent \SI{50}{\pico\second} Langevin NVT simulations at \SI{300}{\kelvin} with friction coefficient $\gamma\!=\!\SI{10}{\per\pico\second}$ and timestep $\Delta t\!=\!\SI{0.25}{\femto\second}$, initialized from 10 distinct test-set conformations and driven by the trained students. The thermostat coupling and timestep are tightened relative to the $\gamma\!=\!\SI{1}{\per\pico\second}$, $\Delta t\!=\!\SI{0.5}{\femto\second}$ commonly used for MD17 small molecules; stronger coupling has been recommended for ML force-field MD on biomolecules to suppress integration instability~\cite{bizzi2025stability,fu2023forces} and is consistent with practice in the MACE and NequIP example simulations ($\gamma\!\approx\!\SI{5}{\per\pico\second}$). Figure~\ref{fig:chignolin_case_study}(c) shows the per-frame C$\alpha$ RMSD relative to each trajectory's initial conformation, averaged over 10 trajectories per method (mean $\pm$ std). In the late-simulation phase (last quarter, $37.5$--$\SI{50}{\pico\second}$), EDG++ attains a lower mean C$\alpha$ RMSD than the baseline (\SI{2.32}{\angstrom} vs.\ \SI{2.60}{\angstrom}, \textbf{10.5\% reduction}), indicating that selective distillation preserves conformational fidelity better than the baseline over long simulations.

\textbf{Inference-time OOD detection: a capability that emerges for free.} Beyond defending against negative transfer at training time---the design role for which the reliability estimator was introduced---the same frozen component repurposes at inference as a per-frame OOD detector for MD rollouts. Figure~\ref{fig:chignolin_case_study}(d) evaluates the reliability score on every sampled MD frame from four EDG++ trajectories: as the simulation drifts toward less-sampled conformations, reliability drops monotonically (e.g., trajectory \#7419 declines from $-0.77$ to $-1.80$ within \SI{25}{\pico\second} and 89\% of its frames fall outside the training $\pm 2\sigma$ band), while the trajectory remaining inside the basin (\#1571) keeps its reliability score within the band throughout. This per-frame reliability signal is unavailable from the student alone, requires no additional training, and emerges \emph{for free} from our training-time design---repositioning the reliability estimator as a general-purpose tool for distributional-drift monitoring during long MD rollouts.

\begin{figure*}[!t]
\centering
\includegraphics[width=0.98\textwidth]{imgs_edgpp/experiments/chignolin_case_study/chignolin_case_study_final.pdf}
\caption{Chignolin case study: a smoking-gun stress test for negative transfer and an inference-time OOD demo. Chignolin is a 166-atom mini-protein substantially out of the ED teacher's training distribution. \textbf{(a) Subject.} Folded native structure (CLN025, PDB 5AWL); the four anchor residues (D3, E5, G7, T8) define the reaction coordinates used in (b). \textbf{(b) Negative transfer, observed.} 2D density landscape of the 9{,}543 DFT-labeled conformations on $d_{D3\text{-}G7}$ vs.\ $d_{E5\text{-}T8}$, with six held-out conformations (paired rescue + amplify across folded/transition/unfolded basins) showing per-conformation energy errors for ViSNet / +EDG / +EDG++; naive EDG degrades the baseline on the rescue cases. \textbf{(c) Gating recovers stability.} C$\alpha$ RMSD over \SI{50}{\pico\second} Langevin NVT (10 trajectories per method, mean $\pm$ std); in the late phase ($37.5$--$\SI{50}{\pico\second}$) EDG++ attains 10.5\% lower mean RMSD than the baseline (\SI{2.32}{\angstrom} vs.\ \SI{2.60}{\angstrom}). \textbf{(d) Inference-time OOD detector (emerges for free).} Per-frame teacher reliability along EDG++ MD trajectories: as the simulation drifts toward less-sampled conformations, reliability drops monotonically---a per-frame OOD signal supplied by the same frozen gating module at no extra training cost.}
\label{fig:chignolin_case_study}
\end{figure*}
% ============================================================================
% END CHIGNOLIN CASE STUDY
% ============================================================================

\FloatBarrier
